\documentclass[a4paper]{article}
\usepackage[spanish]{babel}
\usepackage[utf8]{inputenc}
\usepackage{charter}   % tipografia
\usepackage{graphicx}
%\usepackage{makeidx}
\usepackage{paralist} %itemize inline

%\usepackage{float}
%\usepackage{amsmath, amsthm, amssymb}
%\usepackage{amsfonts}
%\usepackage{sectsty}
%\usepackage{charter}
%\usepackage{wrapfig}
%\usepackage{listings}
%\lstset{language=C}


\usepackage{color} % para snipets de codigo coloreados
\usepackage{fancybox}  % para el sbox de los snipets de codigo

\definecolor{litegrey}{gray}{0.94}

% \newenvironment{sidebar}{%
% 	\begin{Sbox}\begin{minipage}{.85\textwidth}}%
% 	{\end{minipage}\end{Sbox}%
% 		\begin{center}\setlength{\fboxsep}{6pt}%
% 		\shadowbox{\TheSbox}\end{center}}
% \newenvironment{warning}{%
% 	\begin{Sbox}\begin{minipage}{.85\textwidth}\sffamily\lite\small\RaggedRight}%
% 	{\end{minipage}\end{Sbox}%
% 		\begin{center}\setlength{\fboxsep}{6pt}%
% 		\colorbox{litegrey}{\TheSbox}\end{center}}

\newenvironment{codesnippet}{%
	\begin{Sbox}\begin{minipage}{\textwidth}\sffamily\small}%
	{\end{minipage}\end{Sbox}%
		\begin{center}%
		\vspace{-0.4cm}\colorbox{litegrey}{\TheSbox}\end{center}\vspace{0.3cm}}



\usepackage{fancyhdr}
\pagestyle{fancy}

%\renewcommand{\chaptermark}[1]{\markboth{#1}{}}
\renewcommand{\sectionmark}[1]{\markright{\thesection\ - #1}}

\fancyhf{}

\fancyhead[LO]{Sección \rightmark} % \thesection\ 
\fancyfoot[LO]{\small{Nombre Apellido, Nombre Apellido, Nombre Apellido}}
\fancyfoot[RO]{\thepage}
\renewcommand{\headrulewidth}{0.5pt}
\renewcommand{\footrulewidth}{0.5pt}
\setlength{\hoffset}{-0.8in}
\setlength{\textwidth}{16cm}
%\setlength{\hoffset}{-1.1cm}
%\setlength{\textwidth}{16cm}
\setlength{\headsep}{0.5cm}
\setlength{\textheight}{25cm}
\setlength{\voffset}{-0.7in}
\setlength{\headwidth}{\textwidth}
\setlength{\headheight}{13.1pt}

\renewcommand{\baselinestretch}{1.1}  % line spacing


% \setcounter{secnumdepth}{2}
\usepackage{underscore}
\usepackage{caratula}
\usepackage{url}


% ******************************************************** %
%              TEMPLATE DE INFORME ORGA2 v0.1              %
% ******************************************************** %
% ******************************************************** %
%                                                          %
% ALGUNOS PAQUETES REQUERIDOS (EN UBUNTU):                 %
% ========================================
%                                                          %
% texlive-latex-base                                       %
% texlive-latex-recommended                                %
% texlive-fonts-recommended                                %
% texlive-latex-extra?                                     %
% texlive-lang-spanish (en ubuntu 13.10)                   %
% ******************************************************** %



\begin{document}


\thispagestyle{empty}
\materia{Organización del Computador II}
\submateria{Segundo Cuatrimestre de 2014}
\titulo{Trabajo Práctico III}
\subtitulo{System Programming - Zombi defense}
\integrante{Aldasoro Agustina}{86/13}{agusaldasoro@gmail.com}
\integrante{Rey Maximiliano}{XXX/XX}{mail}
\integrante{Tirabasso Ignacio}{XXX/XX}{mail}


\maketitle
\newpage

%\thispagestyle{empty}
%\vfill
%\begin{abstract}
%En el presente trabajo se describe la problemática de ...
%\end{abstract}

%\thispagestyle{empty}
%\vspace{3cm}
%\tableofcontents
\newpage


%\normalsize
\newpage
\section{Ejercio 1}
{\large a)} Armamos los cuatro segmentos de la GDT, llam\'andolos: 
\begin{codesnippet}
\begin{verbatim}
    [GDT_IDX_CODE_0] = (gdt_entry) ; 
    [GDT_IDX_CODE_3] = (gdt_entry) ; 
    [GDT_IDX_DATA_0] = (gdt_entry) ;  
    [GDT_IDX_DATA_3] = (gdt_entry) ; 
\end{verbatim}
\end{codesnippet}
    
A los cuatro les seteamos el mismo \emph{l\'imite}: 0x26EFF = ... , y la misma \emph{base} en 0. El \emph{segment type} var\'ia depende el segmento: CODE_0: 0x0A= ... , CODE_3: 0x0F= ... , DATA_0 y DATA_3: 0x02= ... . El \emph{Descriptor type} va en todos para system, por lo tanto es 0. El \emph{Descriptor privilege level} coincide con el nombre del descriptor (0 para CODE_0 y DATA_0; 3 para CODE_3 y DATA_3). El bit de \emph{Present} va para todos en 1 y los bit de \emph{Available for use by system software} y \emph{l} van para todos en 0. El bit de \emph{Default operation size} va para todos en 1 porque es un c\'odigo de 32bits. El bit de \emph{Granularity} va para todos en 1.\\
    


{\large b)} Se adjunta el c\'odigo necesario para pasar a modo protegido y setear la pila del kernel en
la direcci\'on 0x27000.


\begin{codesnippet}
\begin{verbatim}
    ; Deshabilitar interrupciones
    cli
  
    ; Habilitar A20
    call habilitar_A20
	
    ; Cargar la GDT
    lgdt [GDT_DESC]

    ; Setear el bit PE del registro CR0
    mov eax,cr0
    or eax,1
    mov cr0,eax

    jmp 0x50:modo_protegido
\end{verbatim}
\end{codesnippet}
\begin{codesnippet}
\begin{verbatim}
BITS 32
modo_protegido:

    ; Establecer selectores de segmentos
    xor eax, eax
    mov ax, 0x40
     
    mov es, ax
    mov ds, ax
    mov ss, ax    
    mov gs, ax

    mov ax, 0x60 
    mov fs, ax
    
    ; Establecer la base de la pila
    mov ebp, 0x27000
\end{verbatim}
\end{codesnippet}

{\large c)} Segmento adicional que describe el \'area de la pantalla en memoria que puede
ser utilizado solo por el kernel:
\begin{codesnippet}
\begin{verbatim}
    ; Cambiar modo de video a 80 X 50
    mov ax, 0003h
    int 10h ; set mode 03h
    xor bx, bx
    mov ax, 1112h
    int 10h ; load 8x8 font
\end{verbatim}
\end{codesnippet}


{\large d)} La siguiente es la rutina que se encarga de limpiar la pantalla y pintar el \'area del mapa
un fondo de color verde, junto con las dos barras laterales para cada uno de
los jugadores (rojo y azul).
\begin{codesnippet}
\begin{verbatim}
    ; Inicializar pantalla
    call clear_screen
    call print_map ;código que va a pintar el mapa de los colores deseados
\end{verbatim}
\end{codesnippet}
\begin{codesnippet}
\begin{verbatim}
void clear_screen() {
    int size = VIDEO_COLS * VIDEO_FILS;    
    ca (*p) = (ca (*)) VIDEO; 
    int i = 0;
    ca empty;
    empty.c = 0;
    empty.a = getFormat(C_FG_BLACK, 0, C_BG_BLACK, 0);
    while(i < size) {
        p[i] = empty;
        i++;
    }
}
\end{verbatim}
\end{codesnippet}

\newpage
\section{Ejercio 2}
{\large a)} Inicializamos las entradas 0-19 en la IDT.





{\large b)} Para probar lo programado anteriormente hicimos:

\begin{codesnippet}
\begin{verbatim}
    ; Inicializar la IDT
    call idt_inicializar
    ; Cargar IDT
    lidt [IDT_DESC]

    ; test para que salte la divide by 0 exception (0)
    mov edx,0
    mov ecx,0
    mov eax,3
    div ecx
\end{verbatim}
\end{codesnippet}

\newpage
\section{Ejercio 3}
{\large a)} La siguiente es la rutina que se encarga de limpiar el buffer de video y pintarlo como indica
la figura 9. 

\begin{codesnippet}
\begin{verbatim}
void print_map() {
    int cols = VIDEO_COLS;
    int rows = VIDEO_FILS;

    ca (*screen)[VIDEO_COLS] = (ca (*)[VIDEO_COLS]) VIDEO;

	/* Defino los colores que voy a usar */
    ca red;
    red.c = 0;
    red.a = getFormat(C_FG_RED, 0, C_BG_RED, 0);
    ca blue;
    blue.c = 0;
    blue.a = getFormat(C_FG_BLUE, 0, C_BG_BLUE, 0);
    ca green;
    green.c = 0;
    green.a = getFormat(C_FG_GREEN, 0, C_BG_GREEN, 0);
    ca black;
    black.c = 0;
    black.a = getFormat(C_FG_BLACK, 0, C_BG_BLACK, 0);

    int y,x;

    clear_screen();
    
    for(y = 0; y < rows; y++) {
        for(x = 0; x < cols; x++) {
            if (y >= rows-5) {
                screen[y][x] = black;		//las últimas 5 filas van de negro
            } else if (x == cols-1) {
                screen[y][x] = blue;		//el jugador de la derecha es azul
            } else if (x == 0) {
                screen[y][x] = red;			//el jugador de la izquierda es rojo
            } else {
                screen[y][x] = green;		//el terreno de juego es verde
            }
        }
    }
    /*Armo el cuadrado para el puntaje rojo */
    for(y = rows-5; y < rows; y++) {
        for(x = 35; x < 40; x++) {
            screen[y][x] = red;
        }
    }
    /*Armo el cuadrado para el puntaje azul */
    for(y = rows-5; y < rows; y++) {
        for(x = 40; x < 45; x++) {
            screen[y][x] = blue;
        }
    }
}
\end{verbatim}
\end{codesnippet}


{\large b)} Se muestran las rutinas encargadas de inicializar el directorio y tablas de p\'aginas para el kernel
(mmu_inicializar_dir_kernel). 

\begin{codesnippet}
\begin{verbatim}
void mmu_inicializar_dir_kernel() {
    page_directory *pd = (page_directory *) 0x27000;
    int i;

    /* Creo 1024 entradas en page_directory con todo cero. */
    for (i = 0; i < 1024; i++) {
        pd[i] = (page_directory) {};
    }

    /* Mapeo las primeras 4 (0,1,2,3) entradas del page directory
    con la base y permisos correspondientes. */
    for(i = 0; i < 4; i++) {
        pd[i] = (page_directory) {
            .base = 0x28 + i,
            .rw = 0x1,
            .p = 1,
        };
    }

    page_table* pt = (page_table*) 0x28000;
    for(i = 0; i < 1024; i++) {
        pt[i] = (page_table) {
            .base = i,
            .rw = 0x1,
            .p = 1,
        };
    }
}
\end{verbatim}
\end{codesnippet}



{\large c)} Completar el c\'odigo necesario para activar paginaci\'on.\\

\begin{codesnippet}
\begin{verbatim}
    ; Habilitar paginacion

    mov eax,0x27000
    mov cr3,eax

    mov eax,cr0
    or eax,0x80000000
    mov cr0,eax


    mov eax,0x100000
    mov cr3,eax
\end{verbatim}
\end{codesnippet}

\newpage
\section{Ejercio 4}
{\large a)} Para administrar la memoria en el área libre, tenemos un contador de páginas utilizadas denoninándolo \textit{páginas}. Luego contamos con las funciones \textit{get_page_directory} y \textit{get_page_table} las cuales nos brindan un nuevo page directory o una nueva page table correspondientemente.


\begin{codesnippet}
\begin{verbatim}
page_directory* get_page_directory() {
    page_directory* pd = (page_directory*) PAGES;
    pd += paginas * 4096;
    int i;
    for (i = 0; i < 1024; i++) {
        pd[i] = (page_directory) {};
    }

    pd[0] = (page_directory) {
        .base = 0x28,
        .rw = 0x1,
        .p = 1
    };

    paginas++;

    return pd;
}
\end{verbatim}
\end{codesnippet}


\begin{codesnippet}
\begin{verbatim}
page_table* get_page_table() {
    page_table* pgt = (page_table*) PAGES;
    pgt += paginas * 0x1000;
    int i;
    for (i = 0; i < 1024; i++) {
        pgt[i] = (page_table) {};
    }

    paginas++;

    return pgt
    };
}
\end{verbatim}
\end{codesnippet}
{\large b)} La rutina \textit{mmu_inicializar_dir_zombi} se encarga de inicializar un directorio
de p\'aginas y tablas de p\'aginas para una tarea, respetando la figura 6. Copia
el c\'odigo de la tarea a su \'area asignada, es decir la posici\'on indicada por el jugador dentro
del mapa y mapea dichas p\'aginas a partir de la direcci\'on virtual 0x08000000(128MB).


\begin{codesnippet}
\begin{verbatim}
/* Guerrero = 0, Mago = 1, Clerigo = 2
   player = 0 es A
   player = 1 es B
*/
  page_directory* mmu_inicializar_dir_zombie(unsigned int player, unsigned char class, 
  unsigned int y) {

    page_directory* pd = get_page_directory();
    unsigned int x = (player ? 79 : 2);

    unsigned int offset_x[9] = {0, -1, -1, -1, 0, 0, 1, 1, 1};
    unsigned int offset_y[9] = {0, 0, -1, 1, -1, 1, 0, 1, -1};

    /* En este for se mapean las nueve páginas correspondientes a un jugador */
    int i, _x, _y;
    for(i = 0; i < 9; i++) {
        _x = y + offset_x[i] * (player ? 1 : -1);  \\ x es la posición x dentro del mapa
        _y = x + offset_y[i] * (player ? 1 : -1);  \\ y es la posición y dentro del mapa

        /* La función get_physical_address devuelve la dirección física en memoria
        a la cual le corresponde el par (x, y) pasado por parámetro */      
        mmu_mapear_pagina(0x8000000 + (i*0x1000), pd, get_physical_address(_x, _y), 1, 0);      
    }

    /* Código para copiar la tarea del Zombie */
    int address = address = 0x10000 + (player ? 0 : 1)  * 0x3000 + class * 0x1000;
    i = 0;
    unsigned char *code = (unsigned char *) 0x8000000;
    unsigned char * paddress = (unsigned char*) address;
    while (i++ < 0x1000) {
        code[i] = paddress[i];
    }
    
    return pd;
}
\end{verbatim}
\end{codesnippet}



{\large c)} La rutina \textit{mmu_mapear_p\'agina} permite mapear la p\'agina f\'isica correspondiente a su f\'isica en la direcci\'on virtual utilizando cr3.

\begin{codesnippet}
\begin{verbatim}
void mmu_mapear_pagina(unsigned int virtual, page_directory* pd, unsigned int fisica, 
   unsigned char rw, unsigned char us) {

    unsigned int directory = (virtual >> 22);
    unsigned int table     = (virtual & 0x003FF000) >> 12;

    page_table* pt = (page_table*) (pd[directory].base << 12);

    if (pd[directory].p == 0){
        pd[directory].base = ((unsigned int) get_page_table()) >> 12;
        pd[directory].rw = rw;
	    pd[directory].us = us;
        pd[directory].p = 1;
    }

    pt = (page_table*) (pd[directory].base << 12);
    pt[table].base = fisica >> 12;
    pt[table].rw = rw;
    pt[table].us = us;
    pt[table].p = 1;

    tlbflush();
}
\end{verbatim}
\end{codesnippet}


\textcolor{red}{Esto no lo hicimos}
II- mmu unmapear pagina(unsigned int virtual, unsigned int cr3)
Borra el mapeo creado en la direcci\'on virtual virtual utilizando cr3. \\


\begin{codesnippet}
\begin{verbatim}
Aca va la implementacion de muu_unmapear_pagina
\end{verbatim}
\end{codesnippet}

\newpage
\section{Ejercio 5}

{\large a)} Completamos las entradas necesarias en la IDT para asociar una rutina a la interrupci\'on del
reloj, otra a la interrupci\'on de teclado y por \'ultimo una a la interrupci\'on de software 0x66. Es decir las posiciones 32,33 y 66.


{\large b)} A continuación, la rutina asociada a la interrupci\'on del reloj, para que por cada tick llame a la
funci\'on screen pr\'oximo reloj. La misma se encarga de mostrar cada vez que se llame, la
animaci\'on de un cursor rotando en la esquina inferior derecha de la pantalla. 

\begin{codesnippet}
\begin{verbatim}
;; Rutina de atención del RELOJ

global _isr32
_isr32:
    pushad
    call proximo_reloj 	 ; Ya definida en isr.asm
    call proximo_indice  ;Devuelve el próximo índice en la GDT a ejecutar

    cmp ax,0
    je .nojump

    mov [sched_tarea_selector], ax
    call fin_intr_pic1
    jmp far [sched_tarea_offset]
    jmp .end

.nojump:
    call fin_intr_pic1

.end:
    ; switchear tareas.
    popad
    iret
\end{verbatim}
\end{codesnippet}


{\large c)} Ahora, la rutina asociada a la interrupci\'on de teclado de forma que si se presiona cualquiera
de las teclas a utilizar en el juego, se presenta la misma en la esquina superior derecha de la pantalla. \\

\begin{codesnippet}
\begin{verbatim}
;; Rutina de atención del TECLADO

global _isr33
extern printf
extern print_int
extern handle_keyboard_interrumption
_isr33:
    pushad
    xor eax,eax
    in al, 0x60

    mov dword [esp], eax
    call handle_keyboard_interrumption
    
    mov dword  [esp + 0x], 0
    mov dword [esp + 0xc], 67
    mov dword [esp + 0x8], keyboard_str
    mov dword [esp + 0x4], eax
    call printf    ;función implementada por nosotros

    call fin_intr_pic1
    popad
    iret
\end{verbatim}
\end{codesnippet}


{\large d)} Escribir la rutina asociada a la interrupci\'on 0x66 para que modifique el valor de eax por
0x42. Posteriormente este comportamiento va a ser modificado para atender el servicio del
sistema.

\textcolor{red}{Esto no lo hicimos tampoco!}

\newpage
\section{Ejercio 6}
{\large a)} Definimos tres entradas en la GDT que consideramos necesarias para ser usadas como descriptores
de TSS: una para ser utilizada por la tarea inicial, otra para la tarea actual y una \'ultima para la tarea siguiente.\\

{\large b)} Completamos la entrada de la TSS de la tarea Idle con la informaci\'on de la tarea Idle. La tarea Idle se encuentra en la direcci\'on 0x00016000. La pila se alojar\'a en la misma direcci\'on que la pila del kernel y ser\'a mapeada
con identity mapping. Esta tarea ocupa 1 pagina de 4KB y debe ser mapeada con identity mapping. Adem\'as la misma comparte el mismo CR3 que el kernel. \\

\begin{codesnippet}
\begin{verbatim}
void tss_inicializar() {
    int i = 0;
    while(i < CANT_ZOMBIS) {
        inUseA[i] = 0;
        inUseB[i] = 0;
        i++;
    }
    currentZombieA = 0;
    currentZombieB = 0;

    // inicializar tss_idle
    tss_inicializar_tarea_idle();

    memcpy(&tss_idle, &tss_inicial, sizeof(tss));
    memcpy(&tss_idle, &current_task, sizeof(tss));
    memcpy(&tss_idle, &next_task, sizeof(tss));

    gdt[GDT_INITIAL_TSS].base_31_24 = ((u32) (&tss_inicial) & 0xFF000000) >> 24;
    gdt[GDT_INITIAL_TSS].base_23_16 = ((u32) (&tss_inicial) & 0x00FF0000) >> 16;
    gdt[GDT_INITIAL_TSS].base_0_15  = (u32) (&tss_inicial) & 0x0000FFFF;

    gdt[GDT_CURRENT_TSS].base_31_24 = ((u32) (&current_task) & 0xFF000000) >> 24;
    gdt[GDT_CURRENT_TSS].base_23_16 = ((u32) (&current_task) & 0x00FF0000) >> 16;
    gdt[GDT_CURRENT_TSS].base_0_15  = (u32) (&current_task) & 0x0000FFFF;
}
\end{verbatim}
\end{codesnippet}


\begin{codesnippet}
\begin{verbatim}
void tss_inicializar_tarea_idle() {

    tss_idle = (tss) {};

    tss_idle.eip = 0x00016000;
    tss_idle.cr3 = 0x27000;
    
    tss_idle.ebp = 0x27000;
    tss_idle.esp = 0x27000;

    tss_idle.es = 0x40;
    tss_idle.ds = 0x40;
    tss_idle.ss = 0x40;
    tss_idle.gs = 0x40;
    tss_idle.cs = 0x50;

    tss_idle.eflags = 0x202;
    tss_idle.iomap = 0xffff;
}
\end{verbatim}
\end{codesnippet}

{\large c)} Construir una funci\'on que complete una TSS libre con los datos correspondientes a una
tarea (zombi). El c\'odigo de las tareas se encuentra a partir de la direcci\'on 0x00010000
ocupando una pagina de 4kb cada una seg\'un indica la figura 1. Para la direcci\'on de la
pila se debe utilizar el mismo espacio de la tarea, la misma crecer\'a desde la base de la tarea. Para el mapa de memoria se debe construir uno nuevo utilizando la funci\'on mmu inicializar dir zombi. Adem\'as, tener en cuenta que cada tarea utilizar\'a una pila distinta de nivel 0, para esto se debe pedir una nueva pagina libre a tal fin. \\

\textcolor{red}{Y aca que onda?}\\

{\large d)} C\'odigo necesario para ejecutar la tarea Idle, es decir, saltar intercambiando las
TSS, entre la tarea inicial y la tarea Idle:

\begin{codesnippet}
\begin{verbatim}
idle:
    .loopear:
        inc dword [numero]
        cmp dword [numero], 0x4
        jb .imprimir

    .reset_contador:
        mov dword [numero], 0x0

    .imprimir:
        ; Imprimir 'reloj'
        mov ebx, dword [numero]
        add ebx, message1
        imprimir_texto_mp ebx, 1, 0x0f, 49, 76
        mov ebx, chirimbolo_open
        imprimir_texto_mp ebx, 1, 0x0f, 49, 76-1
        mov ebx, chirimbolo_close
        imprimir_texto_mp ebx, 1, 0x0f, 49, 76+1

    jmp .loopear
\end{verbatim}
\end{codesnippet}

\newpage
\section{Ejercio 7}
{\large a)} Construir una funci\'on para inicializar las estructuras de datos del scheduler.\\

{\large b)} Crear la funci\'on sched pr\'oximo indice() que devuelve el \'indice en la GDT de la pr\'oxima
tarea a ser ejecutada. Construir la rutina de forma devuelva una tarea de cada jugador
por vez seg\'un se explica en la secci\'on 3.2\\

{\large c)} Modificar la rutina de la interrupci\'on 0x66, para que implemente el servicio mover seg\'un
se indica en la secci\'on 3.1.1.\\

{\large d)} Modificar el c\'odigo necesario para que se realice el intercambio de tareas por cada ciclo de
reloj. El intercambio se realizar\'a seg\'un indique la funci\'on sched proximo \'indice().\\

{\large e)} Modificar las rutinas de excepciones del procesador para que desalojen a la tarea que
estaba corriendo y corran la pr\'oxima.\\

{\large f)} Implementar el mecanismo de debugging explicado en la secci\'on 3.4 que indicar\'a en pantalla
la raz\'on del desalojo de una tarea.


\newpage
\section{Ejercio 8}
{\large a)} Crear un conjunto de 3 tareas zombis (Guerrero, Mago y Clerigo). Los mismos deber\'an respetar
las restricciones del trabajo pr\'actico, ya que de no hacerlo no podr\'an ser ejecutados
en el sistema implementado por la c\'atedra.

Deben cumplir:

No ocupar m\'as de 4 kb cada uno (tener en cuenta la pila).

Tener como punto de entrada la direcci\'on cero.

Estar compilado para correr desde la direcci\'on 0x08000000.

Utilizar el unico servicio del sistema (mover).

Explicar en pocas palabras qu\'e estrategia utiliza cada uno de los zombis, o en su conjunto
en t\'erminos de defensa y ataque.\\

{\large b)} Si consideran que sus tareas pueden hacer algo mas que completar el primer item de este
ejercicio, y tienen a un audaz campion que se atreva a enfrentarse en el campo de batalla
zombi, entonces pueden enviar el binario de sus tareas a la lista de docentes indicando
los siguientes datos:

Nombre del campion (Alumno de la materia que se presente como jugador)

Nombre de cada uno de las tareas zombi

Estrategia de alimentaci\'on de los zombis (es decir, como se comer\'an los cerebros de
las otras tareas)

Se realizar\'a una competencia a fin de cuatrimestre con premios en/de chocolate para los
primeros puestos. \\

{\large c)} Pelicula y Video Juego favorito sobre Zombis.


\newpage
\section{Conclusiones y trabajo futuro}


\end{document}






%%%%%%%%%%%%%%%%%%%%%%%%%%%%%%%%%%%%%%%%%%%%%%%%%%%%%%%%%%%%%%%%%%%


\section{Objetivos generales}

El objetivo de este Trabajo Práctico es ...


\section{Contexto}

\begin{figure}
  \begin{center}
	\includegraphics[scale=0.66]{imagenes/logouba.jpg}
	\caption{Descripcion de la figura}
	\label{nombreparareferenciar}
  \end{center}
\end{figure}


\paragraph{\textbf{Titulo del parrafo} } Bla bla bla bla.
Esto se muestra en la figura~\ref{nombreparareferenciar}.



\begin{codesnippet}
\begin{verbatim}

struct Pepe {

    ...

};

\end{verbatim}
\end{codesnippet}


%\section{Enunciado y solucion} 
%\subsection{Filtro cropflip}

Programar el filtro \textit{cropflip} en lenguaje C y luego en ASM haciendo 
uso de las instrucciones vectoriales (\textbf{SSE}).

% ******************************************************************************
\vspace*{0.3cm} \noindent
\textbf{Experimento 1.1 - análisis el código generado}

En este experimento vamos a utilizar la herramienta \verb|objdump| para 
verificar como el compilador de C deja ensamblado el código C.

Ejecutar 
\begin{codesnippet}
\begin{verbatim}
objdump -Mintel -D cropflip_c.o
\end{verbatim}
\end{codesnippet}

¿Cómo es el código generado? 
Indicar
\begin{inparaenum}[\itshape a\upshape)]
    \item Por qué cree que hay otras funciones además de \verb|cropflip_c|
    \item Cómo se manipulan las variables locales
    \item Si le parece que ese código generado podría optimizarse
\end{inparaenum}

% ******************************************************************************
%\newpage
\vspace*{0.3cm} \noindent
\textbf{Experimento 1.2 - optimizaciones del compilador}

Compile el código de C con flags de optimización. Por ejemplo, pasando el flag 
\verb|-O1|\footnote{agregando este flag a \texttt{CCFLAGS64} en el makefile}. 
Indicar
\begin{inparaenum}
    \item Qué optimizaciones observa que realizó el compilador
    \item Qué otros flags de optimización brinda el compilador
    \item Los nombres de tres optimizaciones que realizan los compiladores.
\end{inparaenum}

% ------------------------------------------------------------------------------
% ------------------------------------------------------------------------------

\subsection{Mediciones}

Realizar una medición de performance \emph{rigurosa} es más difícil de lo 
que parece. 
En este experimento deberá realizar distintas mediciones de performance 
para verificar que sean buenas mediciones.

En un sistema ``ideal'' el proceso medido corre solo, sin ninguna 
interferencia de agentes externos. 
Sin embargo, una PC no es un sistema ideal. 
Nuestro proceso corre junto con decenas de otros, tanto de usuarios como 
del sistema operativo que compiten por el uso de la CPU. 
Esto implica que al realizar mediciones aparezcan ``ruidos'' o 
``interferencias'' que distorsionen los resultados.

El primer paso para tener una idea de si la medición es buena o no, 
es tomar varias muestras. 
Es decir, repetir la misma medición varias veces.
Luego de eso, es conveniente descartar los outliers
\footnote{en español, valor atípico: \url{http://es.wikipedia.org/wiki/Valor_atípico}}, 
que son los valores que más se alejan del promedio. 
Con los valores de las mediciones resultantes se puede calcular el promedio 
y también la varianza, que es algo similar el promedio de las distancias al 
promedio\footnote{en realidad, elevadas al cuadrado en vez de tomar el módulo}.

Las fórmulas para calcular el promedio $\mu$ y la varianza $\sigma^2$ son

$$
\mu = \frac{1}{n}\sum_{i=1}^{n} x_i \qquad \sigma^2 = \frac{\displaystyle\sum_{i=1}^{n}(x_i - \mu)^2} {n}
$$

% ******************************************************************************
\newpage
\vspace*{0.3cm} \noindent
\textbf{Experimento 1.3 - calidad de las mediciones}

\begin{enumerate}
    \item Medir el tiempo de ejecución de cropflip 10 veces. 
    \item Implementar un programa en C que no haga más que ciclar 
            infinitamente sumando 1 a una variable. 
            Lanzar este programa tantas veces como \emph{cores lógicos} tenga 
            su procesador. 
            Medir otras 10 veces mientras estos programas corren de fondo.
    \item Calcular el promedio y la varianza en ambos casos.
    \item Consideraremos outliers a los 2 mayores tiempos
     de ejecución de la medicion a) y también a los 2 menores,
     por lo que los descartaremos. Recalcular el promedio y la varianza después de hacer este descarte.
    \item Realizar un gráfico que presente estos dos últimos items.
\end{enumerate}

A partir de aquí todos los experimentos de mediciones deberán hacerse igual 
que en el presente ejercicio: tomando 10 mediciones, luego descartando 
outliers y finalmente calculando promedio y varianza.

% ******************************************************************************
%\newpage
\noindent\textbf{Experimento 1.4 - secuencial vs. vectorial}

En este experimento deberá realizar una medición de las diferencias de 
performance entre las versiones de C y ASM (el primero con -O0, -O1, -O2 y -O3) 
y graficar los resultados.

% ******************************************************************************
\vspace*{0.3cm} \noindent
\textbf{Experimento 1.5 - cpu vs. bus de memoria}

Se desea conocer cual es el mayor limitante a la
performance de este filtro en su versión ASM.

¿Cuál es el factor que limita la performance en este caso?
En caso de que el limitante fuera la intensidad de cómputo, entonces 
podrían agregarse instrucciones que realicen accesos a memoria extra y la
performance casi no debería sufrir. 
La inversa puede aplicarse, si el limitante es la cantidad de accesos a memoria.
\footnote{también podría pasar que estén más bien balanceados y que agregar
cualquier tipo de instrucción afecte sensiblemente la performance}
	
Realizar un experimento, agregando 4, 8 y 16 instrucciones aritméticas 
(por ej \verb|add rax, rbx|) analizando como varía el tiempo de ejecución.
Hacer lo mismo ahora con instrucciones de acceso a memoria, haciendo 
mitad lecturas y mitad escrituras (por ejemplo, agregando dos 
\verb|mov rax, [rsp]| y dos \verb|mov [rsp+8], rax|).\footnote{Notar que en el caso de acceder a \texttt{[rbp]} o \texttt{[rsp+8]} probablemente haya siempre hits en la cache, por lo que la medición no será de buena calidad. Si se le ocurre la manera, realizar accesos a otras direcciones alternativas.}
	
Realizar un único gráfico que compare:
\begin{inparaenum}
    \item La versión original
    \item Las versiones con más instrucciones aritméticas
    \item Las versiones com más accesos a memoria
\end{inparaenum}

Acompañar al gráfico con una tabla que indique los valores graficados.  
  
%\vspace*{0.3cm} \noindent
%\textbf{Experimento 1.6 (\textit{opcional}) - secuencial vs. vectorial (parte II)}
%
%
%Si vemos a los pixeles como una tira muy larga de
%bytes, este filtro en realidad no requiere \emph{casi}
%ningún procesamiento de datos en paralelo. Esto podría
%significar que la velocidad del filtro de C puede
%aumentarse hasta casi alcanzar la del de ASM. ¿ocurre esto?
%	
%Modificar el filtro para que en vez de acceder
%a los bytes de a uno a la vez se accedan como
%tiras de 64 bits y analizar la performance.

% ------------------------------------------------------------------------------
% ------------------------------------------------------------------------------

\subsection*{Filtro \textit{Sierpinski}}

Programar el filtro \textit{Sierpinski} en lenguaje C y en en ASM haciendo 
uso de las instrucciones vectoriales (\textbf{SSE}).

% ******************************************************************************
\vspace*{0.3cm} \noindent
\textbf{Experimento 2.1 - secuencial vs. vectorial}

Analizar cuales son las diferencias de performace entre las versiones de C 
y ASM de este filtro, de igual modo que para el experimento 1.4.

% ******************************************************************************
\vspace*{0.3cm} \noindent
\textbf{Experimento 2.1 - cpu vs. bus de memoria}

¿Cuál es el factor que limita la performance en este filtro?
Repetir el experimento 1.5 para este filtro.

\subsection*{Filtro \textit{Bandas}}

Programar el filtro \textit{Bandas} en lenguaje C y en en ASM haciendo uso de 
las instrucciones vectoriales (\textbf{SSE}).

% ******************************************************************************
\vspace*{0.3cm} \noindent
\textbf{Experimento 3.1 - saltos condicionales}

Se desea conocer que tanto impactan los saltos condicionales en el código 
de filtro Bandas con \verb|-O1| (la versión en C).\\
Para poder medir esto de manera aproximada, remover el código
que detecta a que banda pertenece cada pixel, dejando
sólo una banda.
Por más que la imagen resultante no sea correcta, será posible tomar una
medida aproximada del impacto de los saltos condicionales.
Analizar como varía la performance. 

% ******************************************************************************
\vspace*{0.3cm} \noindent
\textbf{Experimento 3.2 - secuencial vs. vectorial}

Repetir el experimento 1.4 para este filtro.

% ------------------------------------------------------------------------------
% ------------------------------------------------------------------------------

\subsection*{Filtro \textit{Motion Blur}}
Programar el filtro \textit{mblur} en lenguaje C y en ASM haciendo uso de 
las instrucciones \textbf{SSE}.

% ******************************************************************************
\vspace*{0.3cm} \noindent
\textbf{Experimento 4.1}

Repetir el experimento 1.4 para este filtro

